
\subsection{Proof of Lemma~\ref{lem:pnpi-factorization}}\label{proof:pnpi-factorization}
Given the operators $\mcl Q, \mcl{S}, \mcl{R} \in \Pi_4$
\begin{equation*}
\mcl V=\bmat{\mcl Q & \mcl S \\ \mcl S^* & \mcl R}.
\end{equation*}
Since $\mcl V$ is a strict PN-PI operator, there exists $\epsilon>0$ such that
\begin{equation*}
\mcl Q \ge \epsilon I, \qquad \mcl R\le -\epsilon I.
\end{equation*}
In particular, $\mcl Q$ is self-adjoint and coercive. Hence, by the positive
square-root property for coercive self-adjoint PI operators, there
exists a unique positive invertible operator $\mcl S_{11}$ such that
\begin{equation*}
{\mcl Q=\mcl S_{11}^*\mcl S_{11}}, \quad
\mcl S_{12}:=\mcl S_{11}^{-*}\mcl S, \quad
\mcl S_{11}^*\mcl S_{12}=\mcl S.
\end{equation*}
Next define
\begin{equation*}
\mcl W:=\mcl S_{12}^*\mcl S_{12}-\mcl R=\mcl S^*\mcl Q^{-1}\mcl S-\mcl R.
\end{equation*}
Since $\mcl S^*\mcl Q^{-1}\mcl S\ge 0$ and $-\mcl R\ge \epsilon I$, it follows
that
\begin{equation*}
\mcl W\ge \epsilon I.
\end{equation*}
Hence $\mcl W$ is self-adjoint, coercive, and positive. Again by the positive
square-root property, there exists a unique positive invertible operator
$\mcl S_{22}$ such that
\begin{equation*}
\mcl W=\mcl S_{22}^*\mcl S_{22}.
\end{equation*}
Therefore,
\begin{equation*}
\mcl R=\mcl S_{12}^*\mcl S_{12}-\mcl S_{22}^*\mcl S_{22}.
\end{equation*}
Combining these identities gives
\begin{equation*}
\mcl V=\bmat{\mcl S_{11} & \mcl S_{12} \\ 0 & \mcl S_{22}}^*\bmat{I & 0 \\ 0 &
-I}\bmat{\mcl S_{11} & \mcl S_{12} \\ 0 & \mcl S_{22}}.
\end{equation*}
By construction, $\mcl S_{11}>0$ and $\mcl S_{22}>0$.

\subsection{Proof of Lemma~\ref{lem:primal-dual-pnpi}}\label{proof:primal-dual-pnpi}
Given the operators $\mcl Q, \mcl{S}, \mcl{R} \in \Pi_4$, write
\begin{equation*}
\mcl K:=\bmat{-\mcl R & \mcl S^* \\ \mcl S & -\mcl{Q}},
\qquad
\bar{\mcl V}=\mcl K^{-1}.
\end{equation*}
Since \(\mcl V\) is a strict PN-PI operator, there exists \(\epsilon>0\) such
that
\begin{equation*}
\mcl Q\succeq \epsilon I,
\qquad
\mcl R\preceq -\epsilon I.
\end{equation*}
so both \(-\mcl R\) and \(\mcl Q\) are self-adjoint, coercive, and invertible.
Applying the block inverse formula to \(\mcl K\), generalized to the PI-operator
setting from \cite{lu2002inverses}, gives the following expressions. In
addition, \cite[Theorem 2]{peet2018convexsolutionhinftyoptimalcontroller} shows
that the inverse of a coercive, self-adjoint PI operator is again a coercive,
self-adjoint PI operator.
\begin{equation*}
\bar{\mcl Q}
=
\left(\mcl S^*\mcl Q^{-1}\mcl S -\mcl R\right)^{-1}
{\succeq \epsilon I},
\end{equation*}
and
\begin{equation*}
\bar{\mcl R}
=
-\left(\mcl Q+\mcl S(-\mcl R)^{-1}\mcl S^*\right)^{-1}
{\preceq -\epsilon I}.
\end{equation*}
Moreover, by composition algebra $\bar{\mcl S}$ is also a PI operator,
\begin{equation*}
\bar{\mcl S} = {\mcl Q}^{-1} \mcl{S}\bar{\mcl{Q}}
\end{equation*}
Hence \(\bar{\mcl V}\) is a strict PN-PI operator. The converse follows by the
same argument with \(\mcl V\) and \(\bar{\mcl V}\) interchanged, since
\begin{equation*}
\mcl V=\bmat{-\bar{\mcl R} & \bar{\mcl S}^* \\ \bar{\mcl S} & -\bar{\mcl
Q}}^{-1}.
\end{equation*}
Thus \(\mcl V\) is a strict PN-PI operator if and only if \(\bar{\mcl V}\) is a
strict PN-PI operator.

